Considera el sistema homogeni següent, en què $m$ és un nombre real:
\[
\left.\begin{aligned}
x+y+z&=0\\
x+my+2z&=0\\
mx+y+2z&=0
\end{aligned}\right\}
\]

\begin{apartats}

\apartat[1,25]{1}
Discuteix el sistema segons els valors de $m$.

\begin{solucio}
Un sistema homogeni sempre és compatible: $(0,0,0)$ n'és solució, i $\operatorname{rang}A=\operatorname{rang}A^*$
perquè la columna afegida és nul·la.
\[
|A|=\begin{vmatrix}1&1&1\\1&m&2\\m&1&2\end{vmatrix}=2m+2m+1-m^2-2-2=-m^2+4m-3=-(m-1)(m-3).
\]
\begin{itemize}
\item Si $m\ne1$ i $m\ne3$: $\operatorname{rang}A=3$, compatible determinat. L'única solució és
$(0,0,0)$.
\item Si $m=1$ o $m=3$: $|A|=0$ i $\begin{vmatrix}1&1\\1&2\end{vmatrix}=1\ne0$ (files 1 i 2, columnes 1
i 3), de manera que $\operatorname{rang}A=2<3$: compatible indeterminat.
\end{itemize}
\end{solucio}

\begin{tria}{solucions-no-trivials}
\itemtria{resol-m3}{0,75}{1,25}
Resol el sistema per a $m=3$.

\begin{solucio}
Per a $m=3$ n'hi ha prou amb les dues primeres equacions:
$\left.\begin{aligned}x+y+z&=0\\x+3y+2z&=0\end{aligned}\right\}$. Restant-les, $2y+z=0$, $z=-2y$; i
$x=-y-z=y$. Amb $y=\lambda$, les solucions són $(x,y,z)=(\lambda,\ \lambda,\ -2\lambda)$, amb
$\lambda\in\mathbb R$. (Comprovació a la tercera: $3\lambda+\lambda-4\lambda=0$.)
\end{solucio}

\itemtria{amb-z-1}{0,75}{1,25}
Per a quins valors de $m$ el sistema té alguna solució amb $z=1$? Troba-la.

\begin{solucio}
Si $m\ne1$ i $m\ne3$, l'única solució és $(0,0,0)$, amb $z=0$. Cal mirar els dos casos indeterminats.\\
Per a $m=1$, les dues últimes equacions són iguals, $x+y+2z=0$, i restant-hi la primera surt $z=0$:
cap solució té $z=1$.\\
Per a $m=3$, les solucions són $(\lambda,\lambda,-2\lambda)$: $-2\lambda=1$ dona $\lambda=-\frac12$, i la
solució és $\left(-\frac12,-\frac12,1\right)$. Només passa per a $m=3$.
\end{solucio}
\end{tria}

\begin{nomesllarg}
\apartat{0,75}
Discuteix el sistema $\left.\begin{aligned}mx+y&=2\\x+my&=m+1\end{aligned}\right\}$ segons els valors
de $m$.

\begin{solucio}
$|A|=\begin{vmatrix}m&1\\1&m\end{vmatrix}=m^2-1=(m-1)(m+1)$.
\begin{itemize}
\item Si $m\ne\pm1$: $\operatorname{rang}A=\operatorname{rang}A^*=2$, compatible determinat.
\item Si $m=1$: les dues equacions són $x+y=2$: $\operatorname{rang}A=\operatorname{rang}A^*=1<2$,
compatible indeterminat.
\item Si $m=-1$: $\left.\begin{aligned}-x+y&=2\\x-y&=0\end{aligned}\right\}$. $\operatorname{rang}A=1$, i
$\begin{vmatrix}-1&2\\1&0\end{vmatrix}=-2\ne0$ dona $\operatorname{rang}A^*=2$: incompatible.
\end{itemize}
\end{solucio}
\end{nomesllarg}

\end{apartats}
